\chapter{Towards Superintelligence Alignment}
\label{ch:towards-alignment}

\begin{chapterthesis}
\textbf{[STUB]} One paragraph stating the chapter's core claim.
\end{chapterthesis}

\begin{refsection}

\section{Why This Matters}

\textbf{[STUB]} Explain the decision relevance.

\section{Plain-Language Model}

\textbf{[STUB]} Explain the idea without equations.

\section{Formal Model}

\textbf{[STUB]} Introduce variables, assumptions, and equations.

\section{Worked Example}

\textbf{[STUB]} Give one concrete example.

\section{Counterexample or Failure Mode}

\textbf{[STUB]} Give at least one case where the model fails, overreaches, or needs refinement.

\section{What Would Change This View}

\textbf{[STUB]} List evidence, argument, or empirical results that would weaken the chapter.

\section{Summary}

\begin{itemize}
  \item \textbf{[STUB]} Summary point 1.
  \item \textbf{[STUB]} Summary point 2.
  \item \textbf{[STUB]} Summary point 3.
\end{itemize}

\section*{Chapter References}

This chapter draws together the open-risk synthesis of the International AI Safety Report \autocite{iaisr2025}; the documented limits of reinforcement learning from human feedback \autocite{casper2023rlhflimits}; and work on ontological crises in value-bearing systems \autocite{deblanc2011ontological}.

\printbibliography[heading=subbibliography,title={Chapter References}]

\end{refsection}

